\section{总结与延伸}

\subsection{讲者的核心总结}

Ehsan Adeli 在本讲中构建了从最简单分类器到深度学习基石的完整链条。本讲的核心信息可以总结为：

\begin{enumerate}
    \item \textbf{图像分类是计算机视觉的核心}：几乎所有视觉任务（检测、分割、生成等）都建立在分类的基础上
    \item \textbf{数据驱动是正确范式}：从手工规则转向从数据中学习
    \item \textbf{线性分类器是神经网络的基石}：理解 $f = Wx + b$ 是理解整个深度学习的起点
    \item \textbf{损失函数定义了``好''的标准}：没有损失函数就没有优化，没有优化就没有学习
\end{enumerate}

\subsection{全课知识图谱}

本讲建立了一条从最简单到最核心的认知链条：

\begin{center}
\begin{tikzpicture}[node distance=0.9cm, every node/.style={draw, rounded corners, minimum width=4cm, minimum height=0.7cm, font=\small}]
\node (task) {图像分类任务与挑战};
\node (data) [below=of task] {数据驱动方法};
\node (knn) [below=of data] {K 近邻分类器};
\node (hyper) [below=of knn] {超参数调优};
\node (linear) [below=of hyper] {线性分类器};
\node (loss) [below=of linear] {损失函数};
\draw[->, thick] (task) -- (data);
\draw[->, thick] (data) -- (knn);
\draw[->, thick] (knn) -- (hyper);
\draw[->, thick] (hyper) -- (linear);
\draw[->, thick] (linear) -- (loss);
\node[draw=none, right=0.8cm of task, text width=5.5cm, font=\scriptsize] {语义鸿沟、视角/光照/遮挡等变化};
\node[draw=none, right=0.8cm of data, text width=5.5cm, font=\scriptsize] {收集数据 $\rightarrow$ 训练模型 $\rightarrow$ 评估预测};
\node[draw=none, right=0.8cm of knn, text width=5.5cm, font=\scriptsize] {L1/L2 距离、多数投票、$O(N)$ 预测};
\node[draw=none, right=0.8cm of hyper, text width=5.5cm, font=\scriptsize] {训练/验证/测试集、交叉验证};
\node[draw=none, right=0.8cm of linear, text width=5.5cm, font=\scriptsize] {$f=Wx+b$、代数/视觉/几何视角};
\node[draw=none, right=0.8cm of loss, text width=5.5cm, font=\scriptsize] {SVM Hinge Loss、Softmax Cross-Entropy};
\end{tikzpicture}
\end{center}

\subsection{关键 Takeaways}

\begin{importantbox}{五条核心原则}
\begin{enumerate}
    \item \textbf{图像分类的挑战来自语义鸿沟}：计算机看到的是数字张量，而非人类理解的语义内容
    \item \textbf{数据驱动方法取代手工规则}：让算法从大量标注数据中自动学习分类模式
    \item \textbf{超参数调优必须使用验证集}：永远不要在测试集上选择超参数
    \item \textbf{线性分类器是深度学习的基石}：$f = Wx + b$ 是所有神经网络的基本组件
    \item \textbf{损失函数定义优化目标}：SVM 损失关注边距，Softmax 损失关注概率——两者在实践中各有适用场景
\end{enumerate}
\end{importantbox}

\subsection{拓展阅读}

\begin{itemize}
    \item CS231N 课程官网：\url{http://cs231n.stanford.edu/}
    \item CS231N KNN 在线交互式演示：可在课程网站上体验不同 $K$ 值和距离函数的效果
    \item Bishop, \textit{Pattern Recognition and Machine Learning}, Chapter 4 —— 线性分类器的经典参考
    \item Stanford CS229 Machine Learning —— 关于逻辑回归、SVM、交叉验证的更深入讲解
    \item CIFAR-10 数据集：\url{https://www.cs.toronto.edu/~kriz/cifar.html}
    \item Deng et al., ImageNet: A Large-Scale Hierarchical Image Database, CVPR 2009
\end{itemize}

\end{document}
